\documentclass[preview]{standalone}
\usepackage{amsmath}
\usepackage{amssymb}
\usepackage{stellar}
\usepackage{definitions}
\usepackage{bettelini}
\begin{document}
\id{diffeq-continuation}
\genpage
\section{Maximal solutions}

\includesnpt{diffeq-maximal-solution-definition}

\begin{snippetlemma}{diffeq-vector-overlap-uniqueness}{Agreement on an overlap}
    Let \(f\) be continuous and locally Lipschitz in the state, uniformly
    locally in time, on an open \(\Omega\subseteq\realnumbers^{m+1}\).
    Two solutions of \(y'=f(t,y)\) meeting at a point agree on the intersection
    of their interval domains.
\end{snippetlemma}

\begin{snippetproof}{diffeq-vector-overlap-uniqueness-proof}{diffeq-vector-overlap-uniqueness}{Agreement on an overlap}
    On the common interval, the \set \(\{t:u(t)=v(t)\}\) is nonempty,
    relatively closed by continuity, and relatively open by local uniqueness.
    Connectedness makes it the whole interval. This avoids assuming that
    a first point of disagreement or a last point of agreement exists.
\end{snippetproof}

\begin{snippettheorem}{diffeq-maximal-extension-existence}{Maximal extensions}
    Let \(\Omega\subseteq\realnumbers^{m+1}\) be open and \(f\) continuous.
    Every solution of \(y'=f(t,y)\) extends to a maximal solution, whose
    domain is an open interval. If \(f\) is locally Lipschitz in the state,
    uniformly locally in time, the maximal extension is unique.
\end{snippettheorem}

\begin{snippetproof}{diffeq-maximal-extension-existence-proof}{diffeq-maximal-extension-existence}{Maximal extensions}
    Under local uniqueness, collect all extensions \((u_\gamma,I_\gamma)\)
    of the given solution. They agree on overlaps by the preceding lemma.
    Their union is an interval
    \[
        I_*=\bigcup_\gamma I_\gamma
        =(\inf_\gamma\inf I_\gamma,\sup_\gamma\sup I_\gamma).
    \]
    Define \(u_*(t)=u_\gamma(t)\) whenever \(t\in I_\gamma\).
    This is well defined and locally a solution; it is maximal by construction.
    Without uniqueness, order extensions by restriction. Every chain has an
    upper bound obtained by joining its compatible members. Zorn's lemma
    gives a maximal member. A finite endpoint included in its domain would
    allow extension using Peano's theorem there, so a maximal domain is open.
\end{snippetproof}

\section{Continuation at finite endpoints}

\begin{snippetproposition}{diffeq-uniform-continuity-endpoint}{Endpoint limits}
    A uniformly continuous \function \(u\colon(a,b)\fromto\realnumbers^m\)
    has a finite limit at \(b\) if \(b<+\infty\), and at \(a\) if \(a>-\infty\).
\end{snippetproposition}

\begin{snippetproof}{diffeq-uniform-continuity-endpoint-proof}{diffeq-uniform-continuity-endpoint}{Endpoint limits}
    If \(t_n\to b^-\), the sequence \(t_n\) is Cauchy, hence \(u(t_n)\)
    is Cauchy and converges in \(\realnumbers^m\).
    Interleaving any two such sequences proves that their limits agree.
    The sequential criterion gives the endpoint limit; the left endpoint is analogous.
\end{snippetproof}

\begin{snippettheorem}{diffeq-interior-limit-continuation}{Continuation through an interior limit}
    Let \(f\colon\Omega\fromto\realnumbers^m\) be continuous on an open set,
    and let \(u\colon(a,b)\fromto\realnumbers^m\) solve \(u'=f(t,u)\).
    If \(b<+\infty\), \(u(t)\to v\) as \(t\to b^-\), and \((b,v)\in\Omega\),
    then \(u\) extends as a solution beyond \(b\).
\end{snippettheorem}

\begin{snippetproof}{diffeq-interior-limit-continuation-proof}{diffeq-interior-limit-continuation}{Continuation through an interior limit}
    Peano's theorem gives a solution \(\psi\) near \(b\) with \(\psi(b)=v\).
    Join \(u\) for \(t<b\) to \(\psi\) for \(t\geq b\).
    The joined \function is continuous. Moreover
    \[
        u'(t)=f(t,u(t))\longrightarrow f(b,v)=\psi'(b).
    \]
    The endpoint derivative theorem, component by component, shows that its
    left derivative at \(b\) equals its right derivative.
    The join is therefore \(C^1\) and satisfies the equation at \(b\) as well.
\end{snippetproof}

\begin{snippettheorem}{diffeq-escape-compacts}{Escape from compact subsets}
    Let \(\Omega\subseteq\realnumbers^{m+1}\) be open, \(f\) continuous,
    and \(u\colon(a,b)\fromto\realnumbers^m\) a maximal solution.
    At each finite endpoint, its graph eventually leaves every compact
    \(K\subset\Omega\). For example, if \(b<+\infty\), there is \(\delta>0\)
    such that \((t,u(t))\notin K\) for \(b-\delta<t<b\).
    If both endpoints are finite, one can choose a common
    \(0<\delta<(b-a)/2\) for both ends.
\end{snippettheorem}

\begin{snippetproof}{diffeq-escape-compacts-proof}{diffeq-escape-compacts}{Escape from compact subsets}
    Suppose instead that \(t_n\to b^-\) and \((t_n,u(t_n))\in K\).
    Compactness gives, after extraction, \((t_n,u(t_n))\to(b,v)\in K\subset\Omega\).
    Choose a closed cylinder
    \([b-\eta,b+\eta]\times\overline B_r(v)\subset\Omega\);
    let \(M\) bound \(\|f\|\) on it.
    For large \(n\), \(\|u(t_n)-v\|<r/2\) and \(b-t_n<r/(2(M+1))\).
    Until a possible first exit from the ball,
    \(\|u(t)-u(t_n)\|\leq M(t-t_n)<r/2\).
    A first exit before \(b\) is therefore impossible.
    The whole remaining tail stays in that cylinder, so its derivative is bounded.
    It is Lipschitz there and has a finite limit at \(b\); the subsequence
    forces this limit to be \(v\). The interior-limit theorem extends it, a contradiction.
    The left endpoint follows by reversing time.
\end{snippetproof}

\begin{snippetnote}{diffeq-compact-proof-distinction}{Why a subsequence needs an extra argument}
    If an entire solution tail lies in a compact subset of \(\Omega\), boundedness
    of \(f\) immediately gives uniform continuity and an endpoint limit.
    The escape theorem is stronger: it also rules out arbitrarily late returns
    to a compact set. The bounded-speed cylinder argument upgrades those
    returns to a trapped tail, filling the missing step.
\end{snippetnote}

\begin{snippetcorollary}{diffeq-strip-blowup-alternative}{Blow-up alternative on a full strip}
    Let \(f\in C(I\times\realnumbers^m,\realnumbers^m)\), with \(I\) open.
    If a maximal solution has a finite endpoint \(b\in I\), then
    \(\|u(t)\|\to+\infty\) as \(t\to b^-\).
    In particular a bounded solution cannot terminate at an interior time of \(I\).
\end{snippetcorollary}

\begin{snippetproof}{diffeq-strip-blowup-alternative-proof}{diffeq-strip-blowup-alternative}{Blow-up alternative on a full strip}
    If the norm failed to tend to infinity, a sequence of graph points with
    times tending to \(b\) would lie in
    \([b-\eta,b+\eta]\times\overline B_R(0)\subset I\times\realnumbers^m\),
    contradicting escape from compact subsets.
\end{snippetproof}

\begin{snippetexample}{diffeq-oscillating-boundary-example}{No endpoint limit at a domain boundary}
    On \(\Omega=(0,+\infty)\times\realnumbers\), the equation
    \(y'=-\cos(1/x)/x^2\) has solutions \(y=\sin(1/x)+C\).
    They are bounded but have no limit as \(x\to0^+\).
    Their graphs still leave every compact subset of \(\Omega\), since time
    approaches its boundary. A compact subset of an open set has positive
    distance from its complement, so boundary approach cannot remain in it.
\end{snippetexample}
\end{document}
